\subsection{Implicit Function Theorem}

\begin{definition} \label{def:implicit-derivative-blocks} \hlt{Block Splitting of a Linear Map}\\
For $A\in\LL(\R^{n+m},\R^n)$ and $(\h,\mathbf{k})\in\R^n\times\R^m$, define $A_x\h\defeq A(\h,\ze)$ and $A_y\mathbf{k}\defeq A(\ze,\mathbf{k})$.\\
Then $A_x\in\LL(\R^n)$, $A_y\in\LL(\R^m,\R^n)$, and $(\forall\,(\h,\mathbf{k})\in\R^n\times\R^m)\,A(\h,\mathbf{k})=A_x\h+A_y\mathbf{k}$.\\
If $E\subseteq\R^{n+m}$ is open and $\f:E\to\R^n$ is differentiable at $(\x,\y)\in E$, write $D_1\f(\x,\y)\defeq(D\f(\x,\y))_x$ and $D_2\f(\x,\y)\defeq(D\f(\x,\y))_y$.
\end{definition}

\begin{theorem} \label{thm:linear-implicit-function-theorem} \hlt{Linear Implicit Function Theorem}\\
Let $A\in\LL(\R^{n+m},\R^n)$, $A_x$ is invertible. Then $(\forall\,\mathbf{k}\in\R^m)(\exists!\,\h\in\R^n)\,A(\h,\mathbf{k})=\ze$, and this solution is $\h=-(A_x)^{-1}A_y\mathbf{k}$.
\begin{proof}
Fix $\mathbf{k}\in\R^m$, $\h_0\defeq-(A_x)^{-1}A_y\mathbf{k}$. By Definition \ref{def:implicit-derivative-blocks}, $A(\h_0,\mathbf{k})=A_x\h_0+A_y\mathbf{k}=-A_y\mathbf{k}+A_y\mathbf{k}=\ze$, so a solution exists.\\
If $\h\in\R^n$ also satisfies $A(\h,\mathbf{k})=\ze$, then $A_x\h+A_y\mathbf{k}=\ze$, $\h=-(A_x)^{-1}A_y\mathbf{k}=\h_0$; solution exists and is unique.\\
\hlt{Structural Commentary}: Splitting $A$ separates the unknown variable from the parameter.\\
\hlt{Key Insight}: Invertibility of $A_x$ turns the constraint $A(\h,\mathbf{k})=\ze$ into a uniquely solvable linear equation for $\h$.
\end{proof}
\end{theorem}

The implicit theorem below is the nonlinear counterpart of the block equation theorem.\\
The invertible block is the set of variables one can solve for, just as in Theorem \ref{thm:linear-implicit-function-theorem}.

\noindent Figure \ref{fig:implicit-locus-as-graph} depicts the local graph form: writing the solved variable as $\x=\g(\y)$ and marking the parameter value $\b$, the zero locus is parameterized by $\y\mapsto(\g(\y),\y)$ near $(\g(\b),\b)$.
\begin{illustrationfigure}[H]
  \centering
  \begin{tikzpicture}[
  scale=0.95,
  >=Stealth,
  every node/.style={font=\scriptsize},
  axis/.style={->, black!45, line width=0.5pt},
  window/.style={draw=black!55, fill=blue!4, line width=0.6pt},
  locus/.style={blue!70!black, line width=0.9pt},
  guide/.style={black!45, dashed, dash pattern=on 2.5pt off 2.5pt, line width=0.45pt},
  point/.style={circle, fill=black, inner sep=1.2pt}
]
  \draw[axis] (-0.15,0) -- (5.65,0) node[right] {$\mathbf{y}$};
  \draw[axis] (0,-0.15) -- (0,3.55) node[above] {$\mathbf{x}$};

  \draw[window] (2.05,0.92) rectangle (4.80,3.05);

  \coordinate (p) at (3.55,2.05);
  \draw[locus]
    (1.00,0.55) .. controls (1.70,1.10) and (2.35,1.42) .. (p)
    .. controls (4.02,2.38) and (4.70,3.08) .. (5.42,3.18);
  \node[blue!70!black, fill=white, inner sep=1pt] at (4.82,2.82) {$f(x,y)=0$};

  \draw[guide] (p) -- (3.55,0) node[below=2pt] {$b$};
  \draw[guide] (p) -- (0,2.05) node[left=2pt] {$g(b)$};
  \node[point] at (p) {};

  \node[above left=3pt] at (p) {$(g(b),b)$};
  \node[rotate=28, fill=white, inner sep=1pt] at (2.98,1.80) {$x=g(y)$};
\end{tikzpicture}

  \caption{Inside the shaded product neighborhood, the blue zero locus is the graph $\x=\g(\y)$ over the parameter neighborhood.}
  \label{fig:implicit-locus-as-graph}
\end{illustrationfigure}

\begin{theorem} \label{thm:implicit-function-theorem-rn-rm} \hlt{Implicit Function Theorem}\\
Let $E\subseteq\R^{n+m}$ be open, $\f\in C^1(E,\R^n)$, and $(\a,\b)\in E$ satisfy $\f(\a,\b)=\ze$.\\
If $D_1\f(\a,\b)$ is invertible, open neighborhoods $U\ni\a$ in $\R^n$ and $V\ni\b$ in $\R^m$ and a map $\g:V\to U$ can be chosen such that:
\begin{enumerate}[label=(\roman*)]
\setlength{\itemsep}{0pt}
\item $\{(\x,\y)\in U\times V:\f(\x,\y)=\ze\}=\{(\g(\y),\y):\y\in V\}$.\\
The equality $\g(\b)=\a$ holds, and $\g$ is the unique map $V\to U$ whose graph equals this zero set.
\item $\g\in C^1(V,\R^n)$ and $(\forall\,\y\in V)\,D\g(\y)=-\bigl(D_1\f(\g(\y),\y)\bigr)^{-1}D_2\f(\g(\y),\y)$. At $\y=\b$, $D\g(\b)=-\bigl(D_1\f(\a,\b)\bigr)^{-1}D_2\f(\a,\b)$.
\end{enumerate}
\begin{proof}
Set $\Phi:E\to\R^{n+m}$ by $\Phi(\x,\y)\defeq(\f(\x,\y),\y)$, and set $A\defeq D\f(\a,\b)$.\\
Then $\Phi\in C^1(E,\R^{n+m})$ and $(\forall\,(\h,\mathbf{k})\in\R^n\times\R^m)\,D\Phi(\a,\b)(\h,\mathbf{k})=(A_x\h+A_y\mathbf{k},\mathbf{k})$ by Definition \ref{def:implicit-derivative-blocks}.\\
For $(\h,\mathbf{k})\in\R^n\times\R^m$ satisfying $D\Phi(\a,\b)(\h,\mathbf{k})=\ze$, one has $\mathbf{k}=\ze$ and $A_x\h=\ze$, hence $\h=\ze$ by invertibility of $A_x=D_1\f(\a,\b)$.\\
Thus $\ker D\Phi(\a,\b)=\{\ze\}$, so $D\Phi(\a,\b)$ is invertible by Theorem \ref{thm:finite-dimensional-operator-invertible-equivalences}.\\
Open neighborhoods $W\subseteq E$ of $(\a,\b)$, $Z\subseteq\R^{n+m}$ of $(\ze,\b)$ exist such that $\Phi|_W:W\to Z$ is $C^1$ diffeomorphism by Theorem \ref{thm:inverse-function-theorem-rn}.\\
The map $D_1\f$ is continuous by $\f\in C^1(E,\R^n)$ and Definition \ref{def:implicit-derivative-blocks}; the invertible operators form an open set by Theorem \ref{thm:invertible-operators-open-and-inversion-continuous}(ii).\\
Hence choose open neighborhoods $U\subseteq\R^n$ of $\a$ and $V_0\subseteq\R^m$ of $\b$ such that $U\times V_0\subseteq W$ and $(\forall\,(\x,\y)\in U\times V_0)\,D_1\f(\x,\y)$ is invertible.\\
Choose an open neighborhood $V_1\subseteq\R^m$ of $\b$ with $\{\ze\}\times V_1\subseteq Z$ by openness of $Z$. Set $\mathbf q:V_1\to W$ by $\mathbf q(\y)\defeq(\Phi|_W)^{-1}(\ze,\y)$.\\
Then $\mathbf q\in C^1(V_1,\R^{n+m})$ and $\mathbf q(\b)=(\a,\b)$. By continuity of $\mathbf q$ at $\b$, shrink $V_1$ so that $V_1\subseteq V_0$ and $\mathbf q(V_1)\subseteq U\times V_0$.\\
Set $V\defeq V_1$, and let $\g:V\to U$ be the first component of $\mathbf q$.\\
The identity $(\forall\,\y\in V)\,\Phi(\mathbf q(\y))=(\ze,\y)$ implies $(\forall\,\y\in V)\,\mathbf q(\y)=(\g(\y),\y)$; hence $\mathbf q(V)\subseteq U\times V$.
\begin{enumerate}[label=(\roman*)]
\setlength{\itemsep}{0pt}
\item For $(\x,\y)\in U\times V$, $\f(\x,\y)=\ze\Longleftrightarrow\Phi(\x,\y)=(\ze,\y)\Longleftrightarrow(\x,\y)=\mathbf q(\y)\Longleftrightarrow\x=\g(\y)$.\\
The equality $\g(\b)=\a$ follows from $\mathbf q(\b)=(\a,\b)$. If the graph of $\mathbf h:V\to U$ equals the same zero set, then $(\forall\,\y\in V)\,\mathbf h(\y)=\g(\y)$ by the proved equivalence with $\x=\g(\y)$; hence $\mathbf h=\g$.
\item The map $\g$ is $C^1$ as the first component of $\mathbf q$. Fix $\y\in V$ and set $B\defeq D\f(\g(\y),\y)$.\\
Since $(\g(\y),\y)\in U\times V$, the block $B_x=D_1\f(\g(\y),\y)$ is invertible.\\
The identities $(\forall\,\mathbf z\in V)\,\f(\mathbf q(\mathbf z))=\ze$ and $(\forall\,\mathbf{k}\in\R^m)\,D\mathbf q(\y)\mathbf{k}=(D\g(\y)\mathbf{k},\mathbf{k})$ hold by $\Phi(\mathbf q(\mathbf z))=(\ze,\mathbf z)$ and $\mathbf q(\mathbf z)=(\g(\mathbf z),\mathbf z)$.\\
Hence $(\forall\,\mathbf{k}\in\R^m)\,B_xD\g(\y)\mathbf{k}+B_y\mathbf{k}=\ze$ by Theorem \ref{thm:chain-rule-several-variables} and Definition \ref{def:implicit-derivative-blocks}.\\
Therefore $D\g(\y)=-(B_x)^{-1}B_y=-\bigl(D_1\f(\g(\y),\y)\bigr)^{-1}D_2\f(\g(\y),\y)$ by Theorem \ref{thm:linear-implicit-function-theorem} for $B$.\\
At $\y=\b$, this identity becomes $D\g(\b)=-\bigl(D_1\f(\a,\b)\bigr)^{-1}D_2\f(\a,\b)$ by $\g(\b)=\a$.
\end{enumerate}
\hlt{Structural Commentary}: Augmenting $\f$ by the parameter coordinate transfers the zero locus to local inversion.\\
\hlt{Key Insight}: The same invertible $x$-block solves both the graph and its tangent equation.
\end{proof}
\end{theorem}
